\chapter{The details that make it work}
\label{ch:engineering}

\begin{objectives}
\begin{itemize}
\item Explain why a search that evaluates one position at a time wastes almost all of a GPU,
  and what virtual loss does about it.
\item Describe what a collision is, and recognise its fingerprints in engine output ---
  including in this book's own data.
\item Say what Dirichlet noise is for, why it must be \emph{off} in an analysis tool, and
  what happens to reproducibility if it is not.
\item Distinguish a sampled value from a proven one, and explain the Morphy search stopping
  at 188 of a requested 6400 nodes.
\item List the conditions under which a Leela search is exactly reproducible.
\end{itemize}
\end{objectives}

\section{The problem: the search is serial, the hardware is parallel}

Chapter~\ref{ch:byhand}'s loop is strictly sequential. Select a leaf, evaluate it, back the
value up, and only then can you select the next leaf --- because the next selection depends
on the numbers the last backup just changed.

That is a disaster for the hardware. A GPU evaluating one position at a time runs at a few
percent of its capacity: the position is tiny, the network is huge, and almost all the time
goes into moving weights around rather than computing. Evaluating 256 positions at once costs
barely more wall-clock time than evaluating one.

\begin{keyidea}
The central engineering tension in every neural MCTS engine: \textbf{the algorithm wants to
evaluate one position at a time, and the hardware wants hundreds}. Everything in this chapter
is a consequence of that sentence.
\end{keyidea}

\subsection{Virtual loss}

The fix is to select several leaves before evaluating any of them. But naive repetition
selects the \emph{same} leaf every time --- nothing changed between selections, so the argmax
is unchanged.

\textbf{Virtual loss} solves this by pretending, temporarily, that a selected leaf came back
with the worst possible result. When a path is selected but not yet evaluated, add to each
edge on it
\[
  N(s,a) \mathrel{+}= 1, \qquad W(s,a) \mathrel{-}= 1 ,
\]
i.e.\ book a fictitious loss. Now that path looks worse, so the next selection goes somewhere
else. When the real evaluation arrives, the fiction is removed and the true value applied.

\begin{worked}{virtual loss in numbers}
An edge has $N = 9$, $W = 7.2$, so $Q = 0.800$. A search thread selects through it:
\[
  N \to 10,\quad W \to 6.2,\quad Q \to 0.620 .
\]
A drop of $0.18$ --- enough to make a rival edge more attractive, so the second thread
explores elsewhere. When the real value $v = +0.9$ comes back, we undo and apply:
\[
  N \to 10, \qquad W \to 7.2 - (-1) + 0.9 \;\text{i.e.}\; 8.1, \qquad Q = 0.810 .
\]
\emph{(Illustrative numbers.)} Note the size of the effect depends on how many visits the
edge already has: on an edge with $N = 500$ a virtual loss moves $Q$ by $0.004$ and changes
nothing. Virtual loss diversifies exactly where the tree is thin, which is where
diversification is wanted.
\end{worked}

\subsection{Collisions}

Sometimes selection still lands on a leaf that is already queued for evaluation. That is a
\textbf{collision}. The engine can either give up that iteration (wasting it) or record it
and move on. Leela counts collision events, caps them, and by default tolerates a good many.

Collisions explain something visible in this book's own data:

\begin{realdata}[title={Real engine output: requested nodes versus tree size}]
Every search in this book was run with collision handling deliberately turned down
(\texttt{--max-collision-events=1 --max-collision-visits=1 --minibatch-size=1 --threads=1})
so that the search would be strictly sequential and reproducible for the hand simulation.
The cost shows up as a gap between the node budget requested and the tree actually built:
\begin{center}
\begin{tabular}{@{}rrr@{}}
\toprule
Position & \texttt{go nodes} & resulting root visits \\
\midrule
K+P endgame & 800  & 620 \\
Initial position & 1600 & 1493 \\
Morphy position & 400  & 182 \\
\bottomrule
\end{tabular}
\end{center}
The Morphy row has a second cause as well --- see §8.5.
\end{realdata}

\begin{pitfall}
If your tool reports "analysed at 800 nodes", check whether it means \emph{requested} or
\emph{achieved}. The two differ, sometimes by a factor of two, and the difference depends on
settings that have nothing to do with the position. Report the tree size, which is what the
statistics in Chapter~\ref{ch:onemachine} actually depend on.
\end{pitfall}

\section{The evaluation cache}

The same position arrives repeatedly --- through transpositions, through repeated visits to
the same leaf, across successive searches in the same game. Leela keeps a cache of network
outputs keyed by position, so a repeat costs a lookup instead of a forward pass.

This has a subtle consequence you have already met. The cache key includes the position
\emph{history} (Chapter~\ref{ch:byhand}, §7.6), because the history is part of the network's
input. Two identical boards reached by different move orders are different cache entries and
can carry different evaluations. That is not a bug in the cache; it is a fact about the
network's input.

\begin{repobox}[title={in this repository}]
The cache is also why "replay the same analysis and check you get the same numbers" is a
meaningful integrity test for a pipeline --- and why it can fail spuriously if the process
restarts between runs. If you use cache-replay as an identity gate, make sure the comparison
is over positions supplied identically (same history), not just the same FENs.
\end{repobox}

\section{Tree reuse between moves}

When you play the move the engine recommended and the opponent replies, the subtree under
that pair of moves is still valid --- it describes exactly the position now on the board. So
the engine re-roots: it discards the rest and keeps that subtree, complete with its visit
counts.

The practical effect is large. A search that "starts" at 800 nodes may inherit 300 from the
previous move, so its first new node is added to a tree that already has structure. It also
means the analysis you see is not a clean function of the position: identical positions get
different trees depending on how the game arrived there.

\begin{pitfall}
For a \emph{training tool} this is a reproducibility hazard. Analysing a position in the flow
of a game (with tree reuse and history) and analysing the same position cold from a FEN can
give different numbers, sometimes different move rankings. Neither is wrong; they answer
slightly different questions. Pick one convention for your diagnosis pipeline and record it
in the output, so that a number can always be re-derived.
\end{pitfall}

\section{Certainty: when a value stops being an estimate}

Everything in Part~\textsc{ii} treated values as samples to be averaged. But some values are
not estimates at all: a checkmate, a stalemate, a threefold repetition, a position in a
tablebase. These are \textbf{terminal} --- known exactly, no averaging required.

Engines propagate certainty separately from the average. If \emph{any} child of a node is a
proven win for the side to move, the node is a proven win --- and no amount of sampling
elsewhere can change that. Leela's option \texttt{StickyEndgames} (on by default) does exactly
this.

This is what you saw at the end of Chapter~\ref{ch:puct}:

\begin{realdata}[title={Real engine output: a search that stops early}]
Morphy position, \texttt{go nodes 6400}:
\begin{center}
\begin{tabular}{@{}lr@{}}
\toprule
Requested nodes & 6400 \\
Root visits when the engine answered & \textbf{188} \\
Best move & Qb8+ ($n = 3$, $Q = 1.0000$) \\
\bottomrule
\end{tabular}
\end{center}
Once 16.Qb8+ Nxb8 17.Rd8\# was proven, there was nothing left to compute: the root is a
proven win, and continuing would change nothing. The engine reported and stopped. Requesting
6400 nodes produced the identical answer to requesting 1600, in identical time.
\end{realdata}

\begin{keyidea}
A $Q$ of exactly $\pm1.000$ on a small visit count is a signature worth learning to read: it
usually means \emph{proof}, not enthusiasm. A $Q$ of $0.997$ on a large visit count is the
opposite --- overwhelming evidence, no proof. Your tool should not display them the same way.
\end{keyidea}

Related is \textbf{smart pruning}: if the second-best move cannot overtake the best within
the remaining budget, stop early. This is why a Leela search sometimes returns instantly on a
position with one obvious move --- and why "nodes searched" is not a reliable measure of how
hard the engine worked.

\section{Dirichlet noise: the setting you must not leave on}

During \emph{self-play training} Leela deliberately corrupts the root policy with random
noise:
\begin{equation}
  \label{eq:dirichlet}
  P'(a\given s) \;=\; (1-\varepsilon)\,P(a \given s) \;+\; \varepsilon\,\eta_a,
  \qquad \boldsymbol{\eta} \sim \text{Dirichlet}(\alpha),
\end{equation}
typically with $\varepsilon \approx 0.25$. In plain terms: replace a quarter of the policy
with a random distribution, so that occasionally a move the network dislikes gets a serious
look.

The purpose is exploration during \emph{learning}, and it is essential: without it, a network
that has decided a move is bad will never generate the games that could teach it otherwise.
This is the training-time answer to the same problem the $U$ term solves at search time.

\begin{pitfall}
Noise must be \textbf{off} for analysis. With it on:
\begin{itemize}
\item the same position analysed twice gives different answers;
\item the "policy" you display is not the network's opinion but a randomised version of it;
\item any metric you compute about policy --- your policy-blindness measure, the divergence
  index of Chapter~\ref{ch:readingtree} --- is contaminated by noise you injected yourself.
\end{itemize}
Leela's defaults have noise off for normal play, but training-derived configurations and
copied command lines sometimes carry it. If your reruns of the same position disagree, this
is the first thing to check --- and the second is thread count.
\end{pitfall}

\section{When is a search reproducible?}

Determinism matters for a diagnosis pipeline: a metric you cannot recompute is a metric you
cannot debug. Leela's search is deterministic when

\begin{enumerate}
\item there is no Dirichlet noise;
\item temperature is zero (see below), so the move is the argmax rather than a sample;
\item there is a single search thread --- with multiple threads, batches are assembled in a
  race-dependent order and the tree differs run to run;
\item the node limit is a \emph{node} limit, not a time limit (a time limit makes the tree
  size depend on machine load);
\item the cache and the tree start in the same state (Chapter~\ref{ch:byhand}'s runs were all
  cold starts).
\end{enumerate}

Every search in this book satisfies all five, which is why the hand simulation could match to
five decimals. In production you will usually give up (3) for speed, and should then treat
small differences between runs as expected rather than as bugs.

\subsection{Temperature}

For the record, since it appears in configurations: when selecting the move to \emph{play}
from the visit counts, one can either take the maximum or sample from
\begin{equation}
  \label{eq:temperature}
  \pi(a) \;=\; \frac{N(s,a)^{1/\tau}}{\sum_b N(s,b)^{1/\tau}} .
\end{equation}
$\tau \to 0$ gives the argmax (all the weight on the most-visited move); $\tau = 1$ gives
sampling in proportion to visits; large $\tau$ approaches uniform. Self-play uses $\tau = 1$
for the opening moves to generate varied games, then drops to $0$. For analysis, $\tau = 0$.

\begin{worked}{what temperature does to a visit distribution}
Visits $N = (532, 451, 230, 100, 64)$ --- the real numbers from the initial position at 1600
nodes. Total $1377$.
\begin{center}
\begin{tabular}{@{}lrrrrr@{}}
\toprule
 & e4 & d4 & Nf3 & c4 & g3 \\
\midrule
$\tau = 1$ (proportional) & 38.6\% & 32.8\% & 16.7\% & 7.3\% & 4.6\% \\
$\tau = 0$ (argmax)       & 100\%  & 0\%    & 0\%    & 0\%   & 0\%  \\
\bottomrule
\end{tabular}
\end{center}
The $\tau=1$ row is the reason self-play games are not all identical, and --- more
importantly for Chapter~\ref{ch:heads} --- it is the exact quantity used as the training
target for the policy head.
\end{worked}

\section{What all this costs you as a user of the engine}

A summary you can act on:

\begin{center}
\begin{tabular}{@{}lp{9.2cm}@{}}
\toprule
Setting & Why you care \\
\midrule
threads $>1$ & faster; search no longer reproducible \\
minibatch size & bigger = faster on GPU, slightly worse per node (more virtual loss distortion) \\
noise on & \emph{never} for analysis; silently corrupts every policy metric \\
time limit & convenient; makes results depend on machine load \\
tree reuse & realistic; makes a position's analysis depend on the game that reached it \\
FEN vs move history & changes the network input (Chapter~\ref{ch:byhand}, §7.6) \\
\texttt{StickyEndgames} & proven results override sampled ones; explains tiny-$n$ best moves \\
\bottomrule
\end{tabular}
\end{center}

\begin{exercises}

\exhead[virtual loss]{ch8:vloss} An edge has $N = 4$, $W = 3.2$. Compute $Q$ before a virtual
loss, during, and after the real value $+0.7$ arrives. Then repeat for $N = 200$,
$W = 160$. What does the comparison tell you about where virtual loss has its effect?

\exhead[why not more]{ch8:vloss2} Why not apply \emph{two} units of virtual loss, or five, to
push threads further apart? Describe the failure that follows, in terms of the quantity the
search is trying to estimate.

\exhead[collisions]{ch8:collisions} In the K+P run, 800 requested nodes produced a tree of
620. Give two distinct explanations consistent with this chapter, and describe an experiment
that would distinguish them.

\exhead[cache and history]{ch8:cache} Two games transpose into the same position by different
move orders. Will Leela's cache return the same evaluation for both? Justify with reference
to §7.6, and say what this implies for a repertoire tool that stores positions as FENs.

\exhead[tree reuse]{ch8:reuse} Your tool analyses move 25 of a game twice: once by stepping
through the game, once by loading the FEN directly. List three mechanisms in this chapter
that could make the two results differ. Which would you eliminate, and which would you simply
document?

\exhead[certainty]{ch8:certainty} Explain, in terms of §8.4, why the Morphy search answered
after 188 nodes when asked for 6400. Then explain why the same engine, asked for 6400 nodes
in the initial position, uses all of them.

\exhead[reading a Q of 1.0]{ch8:qone} You see two lines in your output:
\texttt{N: 5 (Q: 1.00000)} and \texttt{N: 900 (Q: 0.99312)}. Which is the stronger claim?
Write the one-sentence display text your tool should use for each.

\exhead[noise]{ch8:noise} A colleague reports that policy percentages in your tool "wobble by
a few points between runs". Give the two most likely causes from this chapter, and the
diagnostic that separates them. Why is this a correctness bug and not a cosmetic one for the
policy-blindness metric specifically?

\exhead[temperature]{ch8:temp} Using Equation~\eqref{eq:temperature} with the visit counts
from §8.6.1, compute $\pi$ at $\tau = 0.5$ for e4 and g3. Is the distribution sharper or
flatter than $\tau = 1$? State the general rule.

\exhead[an audit]{ch8:audit} Write a five-item checklist your project should run before
trusting a batch of diagnosis output, drawn only from this chapter. For each item, say what
symptom would show up downstream if it were violated.

\end{exercises}
